\chapter{Введение}
\label{ch:intro}

В оптической отрасли стоит задача обеспечить контроль точности геометрии линз, так как механические методы часто не достигают требуемой погрешности в доли длины волны. Без контроля поверхностей линз создание современных оптических систем становится экономически нецелесообразным. Решением может служить метод колец Ньютона, позволяющий бесконтактно измерять оптические параметры с высокой чувствительностью.

% В основе метода лежит интерференция света в тонком воздушном зазоре между линзой и пластиной. Возникающая картина концентрических колец содержит информацию о геометрии поверхности и свойствах света. Для определения радиуса кривизны линзы измеряют диаметры колец при известной длине волны. Также метод позволяет находить разницу длин волн близких спектральных линий по изменяющей четкость системе колец.
При прохождении света узкополосного спектра через линзу, на поверхности перед ней может возникать интерференционная картина в виде колец, называемая кольцами Ньютона. Возникающая картина концентрических колец содержит информацию о геометрии поверхности и свойствах света. Измеряя диаметры колец и зная длину волны, можно определить радиус кривизны линзы. Если на линзу попадает дополнительная близкая спектральная линия света, то будет наблюдаться изменяющая четкость система колец, по которой можно определить разницу длин волн данных спектральных компонент.

% Целью работы стало определение радиуса кривизны исследуемой линзы и разности длин волн двух спектральных компонент ртутной лампы на основе анализа интерференционных картин в виде колец Ньютона.
Целью работы стало изучение колец Ньютона как метод для определения радиуса кривизны линзы и разности длин волн двух спектральных компонент света.
